\section*{Chapter \ref{ch:b1:seq} --- Sequences}

\begin{solution}{exo:b1:seq:1}
$\Bigl|\dfrac{2n+1}{n+3} - 2\Bigr| = \dfrac{5}{n+3}$. Given
$\varepsilon > 0$, take $N > \frac 5\varepsilon - 3$ (Archimedes):
for $n \geq N$, $\frac{5}{n+3} \leq \varepsilon$. Hence the limit is
$2$.

$((-1)^n)$: its subsequences $(u_{2n}) = (1)$ and $(u_{2n+1}) =
(-1)$ converge to different limits, so the sequence diverges
(\cref{prop:b1:seq:subsequences}). (Directly: any candidate $\ell$
fails for $\varepsilon = \frac12$, since consecutive terms are at
distance $2$.)
\end{solution}

\begin{solution}{exo:b1:seq:2}
Dividing by $n^2$: $\dfrac{1 - 3/n + 1/n^2}{2 + 5/n^2} \to
\dfrac12$.

$\sqrt{n+1} - \sqrt n = \dfrac{1}{\sqrt{n+1} + \sqrt n} \to 0$
(conjugate).

$\dfrac{2^n + n^3}{3^n - n^2} = \dfrac{(2/3)^n + n^3/3^n}{1 -
n^2/3^n} \to \dfrac{0 + 0}{1 - 0} = 0$, using $q^n \to 0$ for $\abs q
< 1$ and the polynomial-vs-geometric comparison
(\cref{prop:b1:functions:powerrules}).

$5^{1/n} = \eu^{\frac{\ln 5}{n}} \to \eu^0 = 1$.
\end{solution}

\begin{solution}{exo:b1:seq:3}
Bernoulli: $(1+h)^n \geq 1 + nh$ for $h \geq -1$, by induction ---
$(1+h)^{n+1} = (1+h)^n(1+h) \geq (1+nh)(1+h) = 1 + (n+1)h + nh^2
\geq 1 + (n+1)h$.

For $0 < \abs q < 1$: write $\frac{1}{\abs q} = 1 + h$, $h > 0$; then
$\abs{q}^n = \frac{1}{(1+h)^n} \leq \frac{1}{1 + nh} \to 0$, and the
squeeze gives $q^n \to 0$ (the case $q = 0$ is trivial). For $q = 1$:
constant sequence, limit $1$. For $q = -1$: diverges
(\cref{exo:b1:seq:1}). For $\abs q > 1$: $\abs q^n = (1 + h)^n \geq
1 + nh \to +\infty$, so $(q^n)$ is unbounded, hence divergent (to
$+\infty$ if $q > 1$; with alternating signs, no limit, if $q <
-1$).
\end{solution}

\begin{solution}{exo:b1:seq:4}
Fixed point: $\ell = \frac{\ell + 3}{2}$ gives $\ell = 3$. Then
\[
v_{n+1} = u_{n+1} - 3 = \frac{u_n + 3}{2} - 3 = \frac{u_n - 3}{2} =
\frac{v_n}{2}:
\]
$(v_n)$ is geometric with ratio $\frac12$, $v_0 = -3$. So $u_n = 3 -
\frac{3}{2^n} \to 3$.
\end{solution}

\begin{solution}{exo:b1:seq:5}
$H_{2n} - H_n = \sum_{k=n+1}^{2n} \frac 1k \geq n \cdot
\frac{1}{2n} = \frac12$ (each of the $n$ terms is $\geq
\frac{1}{2n}$). If $(H_n)$ were Cauchy, taking $\varepsilon =
\frac13$ would force $\abs{H_{2n} - H_n} \leq \frac13$ for large
$n$: contradiction. An increasing non-convergent sequence diverges to
$+\infty$ (\cref{thm:b1:seq:monotone}): $H_n \to +\infty$.
\end{solution}

\begin{solution}{exo:b1:seq:6}
Let $a = \lim u_{2n}$, $b = \lim u_{2n+1}$, $c = \lim u_{3n}$. The
sequence $(u_{6n})$ is a subsequence both of $(u_{2n})$ and of
$(u_{3n})$: its limit equals $a$ and $c$, so $a = c$. The sequence
$(u_{6n+3})$ is a subsequence of $(u_{2n+1})$ (odd indices) and of
$(u_{3n})$ (indices $6n + 3 = 3(2n+1)$): so $b = c$. Hence $a = b$,
and \cref{prop:b1:seq:subsequences} (evens and odds with equal
limits) gives the convergence of $(u_n)$.
\end{solution}

\begin{solution}{exo:b1:seq:7}
\emph{Stability and bounds:} $I = \intcc{0}{2}$ is stable: for $x \in
I$, $\sqrt{2 + x} \in \intcc{\sqrt 2}{2} \subseteq I$; and $u_0 = 0
\in I$.

\emph{Monotonicity:} $f(x) = \sqrt{2+x}$ is increasing and $u_1 =
\sqrt 2 > u_0$: by induction $(u_n)$ is increasing. Increasing and
bounded above by $2$: it converges
(\cref{thm:b1:seq:monotone}).

\emph{Limit:} $\ell = \sqrt{2 + \ell}$ with $\ell \geq 0$ gives
$\ell^2 - \ell - 2 = 0$, so $\ell = 2$.

\emph{Error bound:} multiplying by the conjugate,
\[
2 - u_{n+1} = 2 - \sqrt{2 + u_n}
= \frac{4 - (2 + u_n)}{2 + \sqrt{2+u_n}}
= \frac{2 - u_n}{2 + \sqrt{2 + u_n}}
\leq \frac{2 - u_n}{3},
\]
since $\sqrt{2 + u_n} \geq \sqrt 2 > 1$. By induction from $2 - u_0
= 2$: $\;0 \leq 2 - u_n \leq \frac{2}{3^n}$.
\end{solution}

\begin{solution}{exo:b1:seq:8}
Suppose $(u_n)$ does not converge to $\ell$: for some $\varepsilon_0
> 0$, infinitely many indices satisfy $\abs{u_n - \ell} >
\varepsilon_0$; they form a subsequence $(u_{\varphi(n)})$. This
subsequence is bounded, so by Bolzano--Weierstrass
(\cref{thm:b1:seq:bw}) it has a convergent sub-subsequence, whose
limit $\ell'$ satisfies $\abs{\ell' - \ell} \geq \varepsilon_0$
(pass the inequality to the limit, \cref{thm:b1:seq:order}). But a
sub-subsequence of $(u_n)$ is a convergent subsequence of $(u_n)$,
so by hypothesis $\ell' = \ell$: contradiction.
\end{solution}

\begin{solution}{exo:b1:seq:9}
Suppose $\eu = \frac pq$, $q \geq 1$. The strict inequalities $a_q <
\eu < a_q + \frac{1}{q\,q!}$ (strict since $(a_n)$ is strictly
increasing and $(b_n)$ strictly decreasing) multiplied by $q!$ give
\[
q!\,a_q \;<\; q!\,\frac pq \;<\; q!\,a_q + \frac 1q \leq q!\,a_q + 1.
\]
Now $N = q!\,a_q = \sum_{k=0}^{q} \frac{q!}{k!}$ is an integer (each
$\frac{q!}{k!}$ is a product of integers for $k \leq q$), and so is
$q!\,\frac pq = (q-1)!\,p$. The display thus places the integer
$(q-1)!\,p$ strictly between $N$ and $N + \frac 1q \leq N + 1$: an
integer strictly inside $\intoo{N}{N+1}$, which is impossible. Hence
$\eu \notin \Q$.
\end{solution}

\begin{solution}{exo:b1:seq:10}
\begin{enumerate}
  \item Let $\varepsilon > 0$ and $N$ with $\abs{u_k - \ell} \leq
        \frac{\varepsilon}{2}$ for $k > N$. For $n > N$:
        \[
        \abs{c_n - \ell}
        = \Bigl|\frac{\sum_{k=1}^{n}(u_k - \ell)}{n}\Bigr|
        \leq \frac{\sum_{k=1}^{N} \abs{u_k - \ell}}{n}
        + \frac{n - N}{n}\cdot\frac{\varepsilon}{2}
        \leq \frac{C}{n} + \frac{\varepsilon}{2},
        \]
        where $C = \sum_{k=1}^N \abs{u_k - \ell}$ is fixed. For $n$
        large, $\frac Cn \leq \frac\varepsilon2$: then $\abs{c_n -
        \ell} \leq \varepsilon$.
  \item $u_n = (-1)^n$: diverges, yet $c_n \to 0$ (partial sums
        bounded by $1$, divided by $n$).
  \item Apply (1) to the sequence $v_n = u_{n+1} - u_n \to \ell$: its
        Ces\`aro mean is $\frac{u_{n+1} - u_1}{n} \to \ell$
        (telescoping), and $\frac{u_{n+1}}{n} = \frac{u_{n+1} -
        u_1}{n} + \frac{u_1}{n} \to \ell$; renormalizing indices
        ($\frac{u_n}{n} = \frac{u_n}{n-1}\cdot\frac{n-1}{n}$) gives
        $\frac{u_n}{n} \to \ell$.
\end{enumerate}
\end{solution}

\begin{solution}{exo:b1:seq:11}
Let $L = \inf_{n \geq 1} \frac{u_n}{n} \geq 0$, and $\varepsilon >
0$. Choose $m$ with $\frac{u_m}{m} \leq L + \varepsilon$. Every $n$
writes $n = qm + r$, $0 \leq r < m$; subadditivity (iterated) gives
$u_n \leq q\,u_m + u_r$, so
\[
\frac{u_n}{n} \leq \frac{q m}{n}\cdot\frac{u_m}{m} + \frac{u_r}{n}
\leq \frac{u_m}{m} + \frac{\max(u_0, \dots, u_{m-1})}{n}
\leq L + \varepsilon + \frac{C_m}{n},
\]
using $qm \leq n$. For $n$ large, $\frac{C_m}{n} \leq \varepsilon$:
thus $L \leq \frac{u_n}{n} \leq L + 2\varepsilon$ for all large $n$,
which is the convergence to $L$.
\end{solution}

\begin{solution}{exo:b1:seq:12}
The subgroup $G = \Z + 2\pi\Z$ of $(\R, +)$ is dense: it is not
$\alpha\Z$, since $1 = p\alpha$, $2\pi = q\alpha$ would make $2\pi =
\frac qp$ rational --- and $\pi \notin \Q$ (admitted here; a proof is
sketched in \cref{ch:b1:integration}). By \cref{exo:b1:reals:9}, $G$
is dense in $\R$.

Let now $y \in \intcc{-1}{1}$ and $\theta = \arcsin y$. By density,
for every $\varepsilon > 0$ there are $n \in \Z$, $k \in \Z$ with
$\abs{(n + 2\pi k) - \theta} \leq \varepsilon$, i.e.\ $n$ is within
$\varepsilon$ of $\theta - 2\pi k$; then, $\sin$ being
$2\pi$-periodic and $1$-Lipschitz ($\abs{\sin a - \sin b} \leq
\abs{a - b}$, from the mean value inequality of
\cref{ch:b1:derivative}),
\[
\abs{\sin n - y} = \abs{\sin(n + 2\pi k) - \sin\theta}
\leq \abs{n + 2\pi k - \theta} \leq \varepsilon .
\]
One detail: $n$ ranges over $\Z$, but $\sin(-n) = -\sin n$ and $y$
was arbitrary in $\intcc{-1}{1}$, so nonnegative indices suffice
(replace $(n, y)$ by $(-n, -y)$ if needed). Hence $\{\sin n : n \in
\N\}$ is dense in $\intcc{-1}{1}$; a dense-in-a-segment sequence has
subsequences approaching distinct values, so it diverges.
\end{solution}

\begin{solution}{pb:b1:seq:1}
\textbf{1.} $\dfrac{n(n+1)/2}{n^2} = \dfrac{1 + 1/n}{2} \to
\dfrac12$, and $\dfrac{n(n+1)(2n+1)/6}{n^3} = \dfrac{(1 +
1/n)(2 + 1/n)}{6} \to \dfrac13$.

\textbf{2.} Upper: each of the $n$ terms is $\leq \sqrt n$, so
$T_n \leq n\sqrt n$. Lower: the terms with $k > \frac n2$ number
at least $\frac n2$, and each is $\geq \sqrt{n/2}$:
\[
T_n \geq \frac n2 \sqrt{\frac n2} = \frac{n^{3/2}}{2\sqrt 2} .
\]

\textbf{3.} For $k \geq N$, since $b_{k+1} - b_k > 0$:
$m(b_{k+1} - b_k) \leq a_{k+1} - a_k \leq M(b_{k+1} - b_k)$.
Summing for $k = N, \dots, n - 1$, both sides telescope:
\[
m\,(b_n - b_N) \leq a_n - a_N \leq M\,(b_n - b_N),
\]
and dividing by $b_n - b_N > 0$ gives the claim.

\textbf{4.} By definition of the limit there is $N$ with $\ell -
\varepsilon \leq \frac{a_{k+1} - a_k}{b_{k+1} - b_k} \leq \ell +
\varepsilon$ for all $k \geq N$; question 3 with $m = \ell -
\varepsilon$, $M = \ell + \varepsilon$ transfers the bracketing
to $\frac{a_n - a_N}{b_n - b_N}$.

\textbf{5.} Expanding the right side of the identity:
\[
\frac{a_N - \ell b_N}{b_n} + \frac{a_n - a_N}{b_n} -
\ell\,\frac{b_n - b_N}{b_n}
= \frac{a_n - \ell b_n}{b_n} = \frac{a_n}{b_n} - \ell .
\]
By question 4 the second factor of the product is bounded by
$\varepsilon$ in absolute value, and $0 < 1 - \frac{b_N}{b_n}
\leq 1$ for large $n$, so
\[
\Bigl|\frac{a_n}{b_n} - \ell\Bigr|
\leq \frac{\abs{a_N - \ell b_N}}{b_n} + \varepsilon
\leq 2\varepsilon
\]
as soon as $b_n \geq \frac{\abs{a_N - \ell b_N}}{\varepsilon}$,
which happens eventually since $b_n \to +\infty$. Hence
$\frac{a_n}{b_n} \to \ell$: the Ces\`aro--Stolz theorem.

\textbf{6.} Given $M$, choose $N$ with $\frac{a_{k+1} -
a_k}{b_{k+1} - b_k} \geq M$ for $k \geq N$; the lower half of
question 3 gives $a_n - a_N \geq M(b_n - b_N)$, so
\[
\frac{a_n}{b_n} \geq \frac{a_N}{b_n} +
M\Bigl(1 - \frac{b_N}{b_n}\Bigr) \longrightarrow M .
\]
Beyond some rank, $\frac{a_n}{b_n} \geq M - 1$; as $M$ was
arbitrary, $\frac{a_n}{b_n} \to +\infty$.

\textbf{7.} With $b_n = n$ and $a_n = u_1 + \dots + u_n$: the
increment quotient is $u_{n+1} \to \ell$, so the Ces\`aro mean
$\frac{a_n}{n}$ tends to $\ell$: part (1) of
\cref{exo:b1:seq:10}. With $a_n = u_n$: the increment quotient
is $u_{n+1} - u_n$, giving part (3). Converse: $a_n = (-1)^n$,
$b_n = n$ has $\frac{a_n}{b_n} \to 0$, yet $a_{n+1} - a_n = \pm
2$ alternates: the increment quotient has no limit.

\textbf{8.} Conjugation:
\[
\bigl((1+h)^{3/2} - 1\bigr)\bigl((1+h)^{3/2} + 1\bigr)
= (1+h)^3 - 1 = 3h + 3h^2 + h^3 .
\]
For $h = \frac1n$: $n\bigl((1 + \frac1n)^{3/2} - 1\bigr) =
\frac{3 + 3/n + 1/n^2}{(1 + 1/n)^{3/2} + 1}$, and $1 \leq (1 +
\frac1n)^{3/2} \leq (1 + \frac1n)^2 \to 1$ (squeeze), so the
denominator tends to $2$ and the expression to $\frac32$. Now
Stolz with $a_n = T_n$, $b_n = n^{3/2}$ (strictly increasing,
$\to +\infty$):
\[
\frac{\sqrt{n+1}}{(n+1)^{3/2} - n^{3/2}}
= \frac{\sqrt{n+1}}{\sqrt n} \cdot
\frac{1}{n\bigl((1 + \frac1n)^{3/2} - 1\bigr)}
\longrightarrow 1 \cdot \frac{2}{3},
\]
hence $T_n \sim \frac23\,n^{3/2}$. (The bracketing of question 2
had trapped the constant in $\intcc{0.35}{1}$; Stolz pins it.)

\textbf{9.} (G1) at $u = \ln(1+t)$: $1 + t = \eu^{\ln(1+t)} \geq
1 + \ln(1+t)$, so $\ln(1+t) \leq t$. (G1) at $u =
-\frac{t}{1+t}$: $\eu^{-t/(1+t)} \geq 1 - \frac{t}{1+t} =
\frac{1}{1+t} > 0$; taking $\ln$ (increasing):
$-\frac{t}{1+t} \geq -\ln(1+t)$, i.e.\ $\ln(1+t) \geq
\frac{t}{1+t}$.

\textbf{10.} $\ln$ is strictly increasing
(\cref{prop:b1:functions:expln}), and $\ln(2^k) = k\ln 2$ is
unbounded, so $\ln n \to +\infty$. Increments: with $t =
\frac1n$ in question 9,
\[
\frac 1{n+1} = \frac{1/n}{1 + 1/n} \leq
\ln\Bigl(1 + \frac1n\Bigr) \leq \frac1n
\quad\Longrightarrow\quad
\frac{n}{n+1} \leq \frac{1/(n+1)}{\ln(1 + 1/n)} \leq 1 ,
\]
so the increment quotient $\frac{H_{n+1} - H_n}{\ln(n+1) - \ln
n}$ tends to $1$; Stolz gives $H_n \sim \ln n$.

\textbf{11.} Set $x_n = \frac{u_{n+1}}{u_n} \to L$ and $t_n =
\frac{x_n}{L} - 1 \to 0$. Question 9: $\frac{t_n}{1 + t_n} \leq
\ln(1 + t_n) \leq t_n$, so $\ln x_n - \ln L = \ln(1 + t_n) \to
0$ by squeezing. Ces\`aro (question 7) applied to $(\ln x_k)$:
\[
\frac1n \sum_{k=0}^{n-1} \ln x_k = \frac{\ln u_n - \ln u_0}{n}
\longrightarrow \ln L ,
\]
so $\frac{\ln u_n}{n} \to \ln L$. With $h_n = \frac{\ln u_n}{n}
- \ln L \to 0$: $u_n^{1/n} = L\,\eu^{h_n}$, and (G1) squeezes
$1 + h_n \leq \eu^{h_n} \leq \frac{1}{1 - h_n}$ (for $h_n < 1$),
so $\eu^{h_n} \to 1$ and $u_n^{1/n} \to L$. Application: $u_n =
\binom{2n}{n}$ gives
\[
\frac{u_{n+1}}{u_n} = \frac{(2n+1)(2n+2)}{(n+1)^2}
= \frac{2(2n+1)}{n+1} \longrightarrow 4 ,
\qquad\text{so}\qquad
\binom{2n}{n}^{1/n} \to 4 .
\]

\textbf{12.} For $0 < x \leq 1$, (G2) gives $\sin x \geq x(1 -
\frac{x^2}{6}) \geq \frac{5x}{6} > 0$ and
\[
x - \sin x \geq \frac{x^3}{6} - \frac{x^5}{120}
= x^3\Bigl(\frac16 - \frac{x^2}{120}\Bigr)
\geq \frac{19}{120}\,x^3 > 0 :
\]
thus $0 < \sin x < x$ on $\intoc{0}{1}$. Always $u_1 = \sin u_0
\in \intcc{-1}{1}$. If $u_1 = 0$, then $u_n = 0$ for $n \geq 1$.
If $u_1 \in \intoc{0}{1}$: by induction $0 < u_{n+1} = \sin u_n
< u_n \leq 1$, so $(u_n)_{n\geq1}$ is strictly decreasing and
bounded below by $0$: it converges to some $\ell \in
\intco{0}{1}$ (\cref{thm:b1:seq:monotone}). The product-to-sum
formula and (G2) give $\abs{\sin a - \sin b} = 2\abs{\cos
\frac{a+b}{2}}\,\abs{\sin\frac{a-b}{2}} \leq \abs{a - b}$, so
$u_{n+1} = \sin u_n \to \sin \ell$: $\ell = \sin\ell$. If $\ell
> 0$ then $\sin\ell < \ell$: impossible. So $u_n \to 0$.

\textbf{13.} Dividing (G2) by $u_n > 0$:
\[
1 - \frac{u_n^2}{6} \leq \frac{\sin u_n}{u_n} \leq 1 -
\frac{u_n^2}{6} + \frac{u_n^4}{120} \leq 1 ,
\]
and $u_n \to 0$ squeezes $\frac{\sin u_n}{u_n} \to 1$. Dividing
$x - \sin x$ by $x^3$:
\[
\frac16 - \frac{u_n^2}{120} \leq \frac{u_n - \sin u_n}{u_n^3}
\leq \frac16 \longrightarrow \frac16 .
\]

\textbf{14.} Since $u_{n+1} = \sin u_n$:
\[
w_n = \frac{u_n^2 - \sin^2 u_n}{u_n^2 \sin^2 u_n}
= \frac{(u_n - \sin u_n)(u_n + \sin u_n)}{u_n^2 \sin^2 u_n}
= \frac{u_n - \sin u_n}{u_n^3}\cdot
\frac{u_n + \sin u_n}{u_n}\cdot
\Bigl(\frac{u_n}{\sin u_n}\Bigr)^{2}
\]
(check the powers of $u_n$: $3 + 1 + (-4)$ against the $u_n^2$
downstairs and $u_n^4$ upstairs). By question 13 the three
factors tend to $\frac16$, $2$, $1$: $w_n \to \frac13$.

\textbf{15.} $v_n = \frac{1}{u_n^2}$ has increments $v_{n+1} -
v_n = w_n \to \frac13$, so $\frac{v_n}{n} \to \frac13$ by
\cref{exo:b1:seq:10}~(3): $n u_n^2 \to 3$. Then
\[
\abs{\sqrt n\,u_n - \sqrt 3}
= \frac{\abs{n u_n^2 - 3}}{\sqrt n\,u_n + \sqrt 3}
\leq \frac{\abs{n u_n^2 - 3}}{\sqrt 3} \longrightarrow 0 :
\]
$\sqrt n\,u_n \to \sqrt 3$, i.e.\ $u_n \sim \sqrt{3/n}$.

\textbf{16.} Since $n u_n^2 \to 3$, eventually $2 \leq n u_n^2
\leq 4$, i.e.\ $\sqrt{2/n} \leq u_n \leq 2/\sqrt n$. If $u_n
\leq 10^{-2}$ with $n$ in that range, then $2/n \leq 10^{-4}$:
$n \geq 20\,000$; and $\sqrt{3/n} = 10^{-2}$ at $n = 30\,000$.
Heron's method squares the error at each step --- the digit
count doubles --- because at its fixed point the relevant slope
is $< 1$ in modulus (indeed the iteration is contracting).
Here $\sin' 0 = \cos 0 = 1$: the fixed point is neutral, no
geometric contraction exists, and the decay is governed by the
first nonlinear term $-\frac{x^3}{6}$, hence polynomial. One
step of Heron gains more accuracy than ten thousand steps of
the sine.

\textbf{17.} For arbitrary $u_0$: $u_1 = \sin u_0 \in
\intcc{-1}{1}$. If $u_1 = 0$ the sequence vanishes from rank
$1$. If $u_1 > 0$, Part IV applies verbatim. If $u_1 < 0$, set
$v_n = -u_n$: oddness of $\sin$ gives $v_{n+1} = -\sin u_n =
\sin(-u_n) = \sin v_n$ with $v_1 \in \intoc{0}{1}$, so $v_n \sim
\sqrt{3/n}$, i.e.\ $u_n \sim -\sqrt{3/n}$. In all cases
$\abs{u_n} \sim \sqrt{3/n}$ (or the sequence is eventually $0$):
the fall is universal, only the sign remembers the start.

\textbf{18.} First $\frac{u_{n+1}}{u_n} = 1 - \frac{u_n -
u_{n+1}}{u_n^2}\,u_n \to 1 - a \cdot 0 = 1$. Then
\[
\frac{1}{u_{n+1}} - \frac{1}{u_n}
= \frac{u_n - u_{n+1}}{u_n u_{n+1}}
= \frac{u_n - u_{n+1}}{u_n^2}\cdot\frac{u_n}{u_{n+1}}
\longrightarrow a \cdot 1 = a ,
\]
and \cref{exo:b1:seq:10}~(3) gives $\frac{1}{n u_n} \to a$,
i.e.\ $n u_n \to \frac1a$.

\textbf{19.} $v_n = \frac{1}{u_n}$: $v_{n+1} = \frac{1 +
u_n}{u_n} = v_n + 1$, so $v_n = v_0 + n$ and
\[
u_n = \frac{u_0}{1 + n u_0} , \qquad
n u_n = \frac{n u_0}{1 + n u_0} \longrightarrow 1 .
\]
Lemma check: $u_n - u_{n+1} = \frac{u_n^2}{1 + u_n}$, so
$\frac{u_n - u_{n+1}}{u_n^2} = \frac{1}{1 + u_n} \to 1 = a$, and
question 18 predicts $n u_n \to 1$: exact agreement.

\textbf{20.} Positivity by induction ($\eu^{-u} > 0$);
decreasing since $\eu^{-u_n} < 1$ for $u_n > 0$; hence $u_n \to
\ell \geq 0$ (\cref{thm:b1:seq:monotone}). Continuity bridge:
with $h_n = \ell - u_n \to 0$, $\eu^{-u_n} = \eu^{-\ell}
\eu^{h_n} \to \eu^{-\ell}$ by the (G1) squeeze $1 + h_n \leq
\eu^{h_n} \leq \frac{1}{1 - h_n}$; so $\ell = \ell\,
\eu^{-\ell}$, and $\ell > 0$ would force $\eu^{-\ell} = 1$,
false: $\ell = 0$. For $t > 0$, (G1) gives $\eu^{-t} \geq 1 - t$
and $\eu^{-t} \leq \frac{1}{1 + t}$, so
\[
\frac{1}{1 + t} \leq \frac{1 - \eu^{-t}}{t} \leq 1 .
\]
With $t = u_n$: $\frac{u_n - u_{n+1}}{u_n^2} = \frac{1 -
\eu^{-u_n}}{u_n} \to 1$. Question 18 with $a = 1$: $n u_n \to
1$, so $u_n \sim \frac1n$.

\textbf{21.} As in question 18, $\frac{u_{n+1}}{u_n} = 1 -
\frac{u_n - u_{n+1}}{u_n^3}\,u_n^2 \to 1$. Then
\[
\frac{1}{u_{n+1}^2} - \frac{1}{u_n^2}
= \frac{(u_n - u_{n+1})(u_n + u_{n+1})}{u_n^2 u_{n+1}^2}
= \frac{u_n - u_{n+1}}{u_n^3}\cdot
\frac{u_n + u_{n+1}}{u_n}\cdot
\Bigl(\frac{u_n}{u_{n+1}}\Bigr)^2
\longrightarrow a \cdot 2 \cdot 1 = 2a ,
\]
and \cref{exo:b1:seq:10}~(3) gives $\frac{1}{n u_n^2} \to 2a$:
$n u_n^2 \to \frac{1}{2a}$. For the sine, $a = \frac16$
(question 13): $n u_n^2 \to 3$, exactly Part IV.

\textbf{22.} Increments: $a_{n+1} - a_n = 2^{-n} - 2^{-n-1} =
2^{-n-1} = b_{n+1} - b_n$, so the increment quotient is
constantly $1$. Yet $\frac{a_n}{b_n} = \frac{2 - 2^{-n}}{1 -
2^{-n}} \to 2 \neq 1$. The proof of question 5 breaks at the
boundary term: $\frac{a_N - \ell b_N}{b_n} \to 0$ needed $b_n
\to +\infty$; here (with $\ell = 1$) $a_N - b_N = 1$ and $b_n
\to 1$, so the term tends to $1$ --- precisely the residual gap
$2 - 1$.

\textbf{23.} Stolz with $A_n = \sum_{k=1}^n H_k$ and $B_n = n\ln
n$: $B_{n+1} - B_n = \ln(n+1) + n\ln(1 + \frac1n) > 0$ and $B_n
\to +\infty$. By question 9, $\frac{n}{n+1} \leq n\ln(1 +
\frac1n) \leq 1$, so $B_{n+1} - B_n = \ln(n+1) + \theta_n$ with
$\frac12 \leq \theta_n \leq 1$. Hence
\[
\frac{A_{n+1} - A_n}{B_{n+1} - B_n}
= \frac{H_{n+1}}{\ln(n+1)}\cdot
\frac{1}{1 + \theta_n/\ln(n+1)} \longrightarrow 1 \cdot 1 = 1
\]
(question 10 for the first factor; $\theta_n$ bounded and
$\ln(n+1) \to \infty$ for the second). Stolz concludes:
$\sum_{k=1}^n H_k \sim n\ln n$.

\textbf{24.} If $u_n \to \ell > 0$: as in question 11, $\ln u_n
\to \ln\ell$ (question 9 squeeze on $\ln\frac{u_n}{\ell}$), so
the Ces\`aro means $\frac1n\sum_{k=1}^n \ln u_k \to \ln\ell$,
and the exponential bridge of question 11 gives $(u_1\cdots
u_n)^{1/n} = \exp\bigl(\frac1n\sum\ln u_k\bigr) \to \ell$. If
$u_n \to +\infty$: for any $M$, eventually $u_n \geq \eu^M$, so
$\ln u_n \geq M$: $\ln u_n \to +\infty$; the $+\infty$ Ces\`aro
(question 6, $b_n = n$) gives $\frac1n\sum \ln u_k \to
+\infty$, and (G1) ($\eu^s \geq 1 + s$) sends the geometric mean
to $+\infty$. With $u_n = n$: $(n!)^{1/n} \to +\infty$.

\textbf{25.} (i) Completeness entered only through the monotone
limit theorem, to produce the limits in questions 12 and 20;
the Ces\`aro--Stolz theorem itself is pure
$\varepsilon$-management, valid over $\Q$. (ii) Stolz replaces
$\lim \frac{a_n}{b_n}$ by $\lim$ of the quotient of increments,
exactly as l'Hospital replaces $\lim\frac fg$ by
$\lim\frac{f'}{g'}$ --- the differential twin rests on the mean
value theorem of \cref{ch:b1:derivative}. (iii) Heuristic: if
$f(x) = x - a\,x^{p+1} + o(x^{p+1})$ at the neutral fixed point
$0$, then $\frac{1}{u_{n+1}^p} - \frac{1}{u_n^p} \to pa$ and
$u_n \sim (pan)^{-1/p}$: contact of order $p + 1$ yields decay
$n^{-1/p}$ --- the flatter the graph against the diagonal, the
slower the fall. (iv) The constant: (G2) supplies the cubic
coefficient $\frac16$; the factorization of question 14 doubles
it into the telescope increment $\frac13$; Ces\`aro turns
$\frac{1}{u_n^2}$ into $\frac n3$; inverting and taking roots
delivers $\sqrt{3/n}$.
\end{solution}
